\section{Evidence And Verification Status}
\label{sec:appendix-a-evidence-status}

The table separates the verified construction foundation from the new
mathematical results and the open frontier. The rotation operation's
underlying membership and size behavior is Stainless-verified, but its exact
gap-multiplicity theorem belongs to the mathematical layer.

\begin{center}
\begin{tabular}{>{\raggedright\arraybackslash}p{0.34\textwidth}>{\raggedright\arraybackslash}p{0.28\textwidth}>{\raggedright\arraybackslash}p{0.30\textwidth}}
\hline
\textbf{Claim group} & \textbf{Status} & \textbf{Article location}\\
\hline
Sieve construction and transition mechanics & Stainless-verified foundation & Companion Sieve Sequence article\\
Complete-period, local, weighted, capacity-exhaustion, filter-seven, and copy-block results & Mathematically proved & Sections~3--7, Section~9, and Appendix~C\\
Accepted-boundary discrepancy and residue-collision energy & Exact reductions proved; required estimates open & Section~8\\
\hline
\end{tabular}
\end{center}

External mathematical inputs are confined to two theorems: Bertrand's
postulate supports the exact accepted-local-strikes characterization of
Section~4.3 and the asymptotic corollaries
in Appendix~C, and the Prime Number Theorem (with Mertens estimates) supports
the asymptotic stability-gap discussion of
Section~6.6.

The operational Sieve Sequence construction used by these mathematical
properties is Stainless-verified separately in
\href{https://github.com/thiagomata/prime-numbers/blob/master/articles/chapter6/sieve-sequence.md}{Formal Verification of Sieve
Sequence Stages and Their Transitions}. 
